% theme: microscope_usage
\MajorQuestion{顕微鏡の使い方}
\SetRowSpace{2mm}

\Instruction{顕微鏡の持ち方・置き方・レンズに関する次の問いに答えなさい。}
\QQ{顕微鏡を運ぶとき、一方の手で持つ部分を答えなさい。}{アーム}
\QQ{顕微鏡を運ぶとき、もう一方の手で下から支える部分を答えなさい。}{鏡台}
\QQ{顕微鏡は、どのような場所に置くか。}{水平で、直射日光の当たらない明るい場所}
\QQ{レンズを取り付けるとき、接眼レンズと対物レンズのどちらを先に取り付けるか。}{接眼レンズ}
\QQ{観察を始めるとき、はじめは低倍率と高倍率のどちらにするか。}{低倍率}

\Instruction{顕微鏡による観察の手順に関する次の問いに答えなさい。}
\QQ{視野全体を明るくするために調節する部分を、2つ答えなさい。}{反射鏡、しぼり}
\QQ{プレパラートは、観察したいものが視野のどこにくるようにステージに置くか。}{中央}
\QQ{ピントを合わせるとき、はじめに横から見ながら対物レンズとプレパラートの間をどうするか。}{近づける}
\QQ{ピントを合わせるとき、接眼レンズをのぞきながら、対物レンズとプレパラートの間をどうするか。}{遠ざける}
\QQ{高倍率にするとき、どの部分を回して高倍率の対物レンズにするか。}{レボルバー}

\Instruction{顕微鏡の各部と倍率に関する次の問いに答えなさい。}
\QQ{目を近づけてのぞくレンズを何というか。}{接眼レンズ}
\QQ{観察するものに近い側のレンズを何というか。}{対物レンズ}
\QQ{プレパラートをのせる台を何というか。}{ステージ}
\QQ{ピントを合わせるために回す部分を何というか。}{調節ねじ}
\QQ{接眼レンズが10倍、対物レンズが40倍のとき、顕微鏡の倍率は何倍か。}{400倍}

\Instruction{プレパラートのつくり方に関する次の問いに答えなさい。}
\QQ{観察するものをのせるガラス板を何というか。}{スライドガラス}
\QQ{スライドガラスの上にのせた試料に、スポイトで何を1滴落とすか。}{水}
\QQ{試料の上にかぶせる薄いガラスを何というか。}{カバーガラス}
\QQ{カバーガラスをかぶせるとき、気泡が入らないように支える器具を何というか。}{柄つき針}
\QQ{接眼レンズが15倍、対物レンズが10倍のとき、顕微鏡の倍率は何倍か。}{150倍}
